\numpara{4.11}
\label{II.4.11}
Suppose now that \(G\) is a group scheme over \(S\). By~4.1.4,
\(\Lie(G/S)(S)\) is identified with the group of right-invariant infinitesimal
automorphisms of \(G/S\), that is, by~3.14, with the group of derivations
of \(\mathcal{O}_G\) over \(\OS\) invariant under right translation. Moreover
this identification respects the module structure and is
an{\renewcommand{\thefootnote}{(84)}\nde{One has corrected a sign error in
the original, \cf~\textbf{[DG70]}, \S~II.4, 4.4 and~4.6.}}
anti-isomorphism of Lie algebras, as one sees by arguing as in~4.7.3: let
us set \(\mathcal{O}_I=\OS[t]\) and \(\mathcal{O}_{I'}=\OS[t']\), and let
\(x\in L(I)\) and \(y\in L(I')\). The left translation \(\lambda_x\)
(resp.\ \(\lambda_y\)) is an \(S\)-automorphism of \(G_{I\times I'}\) which
induces the identity on \(G_{I'}\) (resp.\ \(G_I\)) and which corresponds to
an \(\OS\)-automorphism
\[
u=\mathrm{id}+td_x
\qquad\text{resp.}\qquad
v=\mathrm{id}+t'd_y
\]
of \(\mathcal{O}_{G_{I\times I'}}=\mathcal{O}_G[t,t']/(t^2,{t'}^2)\), where
\(d_x\), \(d_y\) are \(\OS\)-derivations of \(\mathcal{O}_G\) invariant under
right translation. As the correspondence between \(S\)-automorphisms of
\(G_{I\times I'}\) and \(\OS\)-automorphisms of \(\mathcal{O}_{G_{I\times I'}}\)
is contravariant, \(\lambda_x\lambda_y\lambda_x^{-1}\lambda_y^{-1}\) corresponds
to \(v^{-1}u^{-1}vu=\mathrm{id}+tt'(d_yd_x-d_xd_y)\). One deduces from this,
by~4.7.3, that the map \(x\mapsto -d_x\) is an isomorphism of Lie algebras
(for more details, see \textbf{[DG70]}, \S~II.4, 4.4 and~4.6). The foregoing
is valid for \(\Lie(G/S)(S')=\Lie(G_{S'}/S')(S')\) for every \(S'\to S\).
One then recovers the classical
definition:{\renewcommand{\thefootnote}{(85)}\nde{One has set out the
original in more detail in what follows; in particular, one has added, for
later references, the numbering 4.11.1 to~4.11.8.}}

\begin{scholium}{4.11.1}
\label{II.4.11.1}
Via the isomorphism \(x\mapsto -d_x\), \(\Lie(G/S)\) is identified with the
functor which, to every \(S'\) over \(S\), associates the
\(\mathrm{O}(S')\)-Lie algebra of the derivations of \(G_{S'}\) with respect
to \(S'\) invariant under right translation.
\end{scholium}

As one already knows, by~4.6.1, that every representable group is good, it
follows from the foregoing:

\begin{corollary}{4.11.2}
\label{II.4.11.2}
Every representable group is very good.
\end{corollary}

Let \(\varepsilon\colon S\to G\) be the unit section of \(G\). Set
\(\omega^1_{G/S}=\varepsilon^*(\Omega^1_{G/S})\) and recall (\cf~3.3) that
\(\Lie(G/S)\) is representable by the vector bundle
\(\mathbb{V}(\omega^1_{G/S})\).

\begin{scholium}{4.11.3}
\label{II.4.11.3}
One has therefore associated functorially, to every \(S\)-group scheme \(G\),
a vector bundle \(\mathbb{L}\mathrm{ie}(G/S)=\mathbb{V}(\omega^1_{G/S})\)
over \(S\), which represents the functor \(\Lie(G/S)\), and hence is equipped
with a structure of \(S\)-scheme in \(\OS\)-Lie algebras. Moreover
(\cf~3.4 and~3.8), this construction commutes with base change and with
finite products.
\end{scholium}

\begin{remarks}{4.11.4}
\label{II.4.11.4}
{\renewcommand{\thefootnote}{(86)}\nde{One has set out what follows in more
detail, taking into account the additions made in Exp.~I, \S~6.8.}}
Let \(\pi\) denote the morphism \(G\to S\).

a) The \(\mathcal{O}_G\)-module \(\Omega^1_{G/S}\) is evidently
\((G\times_S G)\)-equivariant (\cf~I, \S~6) and hence, by~I~6.8.1, one has
\(\Omega^1_{G/S}\simeq\pi^*(\omega^1_{G/S})\). It follows for example that
\(\Omega^1_{G/S}\) is locally free (\resp\ locally free of finite type) if
\(\omega^1_{G/S}\) is, which is in particular the case if \(S\) is the
spectrum of a field (\resp\ if \(S\) is the spectrum of a field and \(G\) is
locally of finite type over \(S\)).

b) Moreover, by~I~6.8.2, \(\omega^1_{G/S}\) is equipped with a canonical
structure of \(G\)-\(\OS\)-module, which induces on
\(\mathbb{V}(\omega^1_{G/S})=\mathbb{L}\mathrm{ie}(G/S)\) the adjoint
operation.{\renewcommand{\thefootnote}{(87)}\nde{For this reason, the linear
operation of \(G\) on \(\omega^1_{G/S}\) could be called the ``pre-adjoint''
operation; in fact, by abuse of language, one will still speak of ``the
adjoint operation'' on \(\omega^1_{G/S}\). Let us point out here a somewhat
more general construction, which will be used in Exp.~III (\cf\ in particular
III~0.8). Suppose given, for every \(Y\to S\), an \(\mathcal{O}_Y\)-module
\(\mathcal{N}(Y)\), and for every \(S\)-morphism \(\varphi\colon Z\to Y\), a
morphism of \(\mathcal{O}_Z\)-modules
\(\mathcal{N}(\varphi)\colon\varphi^*\mathcal{N}(Y)\to\mathcal{N}(Z)\),
functorial in \(\varphi\) (this is the case, for example, if \(\mathcal{J}\)
is a quasi-coherent ideal of \(\OS\) and if one sets
\(\mathcal{N}(Y)=\mathcal{J}\mathcal{O}_Y\), \cf\ Exp.~III); then the
operation of \(G\) on \(\omega^1_{G/S}\) induces a (generalized) ``adjoint
operation'' of \(G\) on the \(S\)-functor in abelian groups which, to every
\(f\colon Y\to S\), associates
\(\Hom_{\mathcal{O}_Y}\bigl(f^*(\omega^1_{G/S}),\mathcal{N}(Y)\bigr)\).}}

c) On the other hand (\cf\ EGA~I, 5.3.11 and~5.4.6), \(\varepsilon\) is an
immersion, and is a closed immersion if \(G\) is separated over \(S\). Thus
\(\omega^1_{G/S}\) is identified with \(\mathcal{I}/\mathcal{I}^2\), where
\(\mathcal{I}\) is the quasi-coherent ideal defining \(\varepsilon(S)\) in an
open \(U\) of \(G\) in which \(\varepsilon(S)\) is closed (if \(G\to S\) is
separated, one may take \(U=G\), and if \(G=\Spec\mathcal{A}(G)\) is affine
over \(S\), \(\mathcal{I}\) is none other than the augmentation ideal of
\(\mathcal{A}(G)\), \ie\ the kernel of
\(\varepsilon^{\natural}\colon\mathcal{A}(G)\to\OS\), \cf~I~4.2).%
{\renewcommand{\thefootnote}{(88)}\nde{One has added the preceding
description of \(\omega^1_{G/S}\), which will be useful later (for example,
in \(\mathrm{VII}_{\mathrm{A}}\), 5.5).}}
\end{remarks}

\begin{remark}{4.11.5}
\label{II.4.11.5}
One can deduce from the isomorphism
\(\Omega^1_{G/S}\simeq\pi^*(\omega^1_{G/S})\) that the \(\OS\)-module
\(\omega^1_{G/S}\) is identified with the sheaf
\(\pi^G_*(\Omega^1_{G/S})\) of ``differentials of \(G\) with respect to \(S\)
invariant on the right'', that is, the sheaf whose sections over an open
\(U\) of \(S\) are the sections of \(\Omega^1_{G/S}\) over \(\pi^{-1}(U)\),
invariant under right translation (\cf~I, 6.8.3, compare also with
\(\mathrm{VII}_{\mathrm{A}}\), 2.4).%
{\renewcommand{\thefootnote}{(89)}\nde{This description of
\(\omega^1_{G/S}\) in terms of invariant differentials will not be used in
the sequel.}}
\end{remark}

\begin{notation}{4.11.6}
\label{II.4.11.6}
One denotes by \(\mathcal{L}\mathrm{ie}(G/S)\) the sheaf of sections of the
vector bundle \(\mathbb{L}\mathrm{ie}(G/S)\to S\); it is the \(\OS\)-module
\((\omega^1_{G/S})^{\vee}
=\mathcal{H}\mathrm{om}_{\OS}(\omega^1_{G/S},\OS)\) dual to
\(\omega^1_{G/S}\) (\cf\ EGA~II~1.7.9). It is equipped with a structure of
\(\OS\)-Lie algebra.

As this construction does not commute with base change (in general), the
Lie algebra structure on this module does not make it possible to reconstruct
the structure of \(S\)-scheme in \(\OS\)-Lie algebras on
\(\mathbb{L}\mathrm{ie}(G/S)\).%
{\renewcommand{\thefootnote}{(90)}\nde{For
\(\mathcal{L}\mathrm{ie}(G/S)\simeq(\omega^1_{G/S})^{\vee}\) does not
necessarily determine \(\omega^1_{G/S}\). For example, if
\(S=\mathbf{A}^1_k\) (\(k\) a field) and if \(G\) is the \(S\)-group whose
fiber at \(s=0\) is \(\mathbf{G}_{\mathrm{a},k}\) and whose other fibers are the trivial
group, then \(\mathcal{L}\mathrm{ie}(G/S)=0\) but
\(\Lie(G_s/k)(k)=k\).}}
\end{notation}

However one has:

\begin{lemma}{4.11.7}
\label{II.4.11.7}
One supposes \(\omega^1_{G/S}\) locally free of finite type (which occurs in
particular if \(G\) is smooth over \(S\) (\cf\ EGA~IV\(_4\), 17.2.4), or if
\(S\) is the spectrum of a field and \(G\) is locally of finite type over
\(S\)). Then
\(\mathcal{L}\mathrm{ie}(G/S)^{\vee}
\simeq(\omega^1_{G/S})^{\vee\vee}
\simeq\omega^1_{G/S}\)
and hence
\[
\mathbb{L}\mathrm{ie}(G/S)
=\mathbb{V}(\omega^1_{G/S})
=\mathbb{V}\bigl(\mathcal{L}\mathrm{ie}(G/S)^{\vee}\bigr)
=\mathbb{W}\bigl(\mathcal{L}\mathrm{ie}(G/S)\bigr)
\]
(the last equality resulting from I~4.6.5.1).
\end{lemma}

Finally, let \(G\to H\) be a monomorphism of group functors. Then
\(\Lie(G/S)\to\Lie(H/S)\) is likewise a monomorphism (\cf~3.7). As every
vector monomorphism of vector bundles is a closed
immersion{\renewcommand{\thefootnote}{(91)}\nde{Let
\(f\colon\mathcal{M}\to\mathcal{N}\) be a morphism of \(\OS\)-modules and
\(\mathcal{P}=\Coker(f)\). If \(\mathbb{V}(\mathcal{N})\to\mathbb{V}(\mathcal{M})\)
is a monomorphism, the surjective morphism
\(\operatorname{Sym}(\mathcal{N})\to\operatorname{Sym}(\mathcal{P})\) factors
through \(\OS\), hence \(\mathcal{P}=0\).}}
one obtains:

\begin{corollary}{4.11.8}
\label{II.4.11.8}
Let \(G\hookrightarrow H\) be a monomorphism of \(S\)-group schemes.
\begin{enumeratei}
\item \(\mathbb{L}\mathrm{ie}(G/S)\to\mathbb{L}\mathrm{ie}(H/S)\) is a closed
immersion, and hence \(\omega^1_{H/S}\to\omega^1_{G/S}\) is an epimorphism.
\item If \(\omega^1_{G/S}\) is locally free of finite type, then the
corresponding morphism
\(\mathcal{L}\mathrm{ie}(G/S)\to\mathcal{L}\mathrm{ie}(H/S)\) is an
isomorphism of \(\mathcal{L}\mathrm{ie}(G/S)\) onto a submodule of
\(\mathcal{L}\mathrm{ie}(H/S)\) which is locally a direct factor.
\end{enumeratei}
\end{corollary}

\sgasection{5}{Computation of some Lie algebras}
\label{II.sec.5}

\par\addvspace{0.55\baselineskip}%
\noindent\textbf{5.1. Examples of Lie algebras: the diagonalizable groups.}\ ---
\label{II.ssec.5.1}
Let \(G=\mathrm{D}_S(M)\) be a diagonalizable group over \(S\) (I~4.4). As the
formation of \(\Lie(G/S)\) commutes with base change, it suffices to carry
out the construction for \(G=\mathrm{D}(M)\). One then has:
\[
G(I_S)
=\Hom_{\mathrm{gr.}}\bigl(M,\Gamma(I_S,\mathcal{O}_{I_S})^{\times}\bigr)
=\Hom_{\mathrm{gr.}}\bigl(M,\Gamma(S,\mathrm{D}\OS)^{\times}\bigr).
\]
Now one has a split exact sequence
\[
1\to\Gamma(S,\OS)\to\Gamma(S,\mathrm{D}\OS)^{\times}
\to\Gamma(S,\OS)^{\times}\to 1,
\]
which gives that \(\Lie(G)(S)\) is identified with
\(\Hom_{\mathrm{gr.}}(M,\mathrm{O}(S))\) equipped with its evident structure
of \(\mathrm{O}(S)\)-module. One therefore obtains, after base change:

\begin{proposition}{5.1}
\label{II.5.1}
One has isomorphisms
\[
\Hom_{S\text{-gr.}}(M_S,\OS)\isomto\Lie(\mathrm{D}_S(M)/S)
\qquad\text{and}\qquad
\mathcal{H}\mathrm{om}_{\mathrm{gr.}}(\widetilde{M}_S,\OS)
\isomto\mathcal{L}\mathrm{ie}(\mathrm{D}_S(M)/S),
\]
(where, in the second isomorphism, \(\widetilde{M}_S\) denotes the constant
sheaf of groups on \(S\) defined by \(M\), and
\(\mathcal{H}\mathrm{om}_{\mathrm{gr.}}\) the sheaf of homomorphisms of
sheaves of groups).
\end{proposition}

\begin{corollary}{5.1.1}
\label{II.5.1.1}
If \(M\) is free of finite type (or, as we shall say later, if
\(\mathrm{D}_S(M)\) is a split torus) then (see~I, 4.6.5 for the definition
of~\(\mathbb{W}\))
\begin{align*}
\mathbb{W}\bigl(\mathcal{L}\mathrm{ie}(\mathrm{D}_S(M)/S)\bigr)
&\isomto\mathbb{L}\mathrm{ie}(\mathrm{D}_S(M)/S),\\
M^{\vee}\otimes_{\ZZ}\OS
&\isomto\mathcal{L}\mathrm{ie}(\mathrm{D}_S(M)/S),
\end{align*}
where \(M^{\vee}\) denotes the dual of the abelian group \(M\). In particular
\[
\OS\isomto\Lie(\mathbf{G}_{\mathrm{m},S}/S)
\qquad\text{and}\qquad
\OS\isomto\mathcal{L}\mathrm{ie}(\mathbf{G}_{\mathrm{m},S}/S).
\]
\end{corollary}

\par\addvspace{0.55\baselineskip}%
\noindent\textbf{5.2. Normalizers and centralizers.}\ ---
\label{II.5.2}
Let us first prove a few lemmas.
Recall (\cf~I~3.1.1) that a sequence \(0\to F'\to F\to F''\to 0\) of
\(\OS\)-modules is said to be \emph{exact} if for every \(S'\to S\) the
sequence \(0\to F'(S')\to F(S')\to F''(S')\to 0\) of
\(\mathrm{O}(S')\)-modules is exact.

Likewise, a sequence \(1\to G'\to G\to G''\to 1\) of \(S\)-groups is said to
be exact if for every \(S'\to S\) the sequence of groups
\(1\to G'(S')\to G(S')\to G''(S')\to 1\) is exact.

\begin{lemma}{5.2.1}
\label{II.5.2.1}
{\renewcommand{\thefootnote}{(92)}\nde{One has set out the statement of the
lemma (and its proof) in more detail; this will be useful in~5.3.1.}}
Let \(1\to G'\to G\to G''\to 1\) be an exact sequence of \(S\)-groups.
\begin{enumeratei}
\item The sequences
\(1\to\mathrm{T}_{G'/S}(\mathcal{M})\to\mathrm{T}_{G/S}(\mathcal{M})
\to\mathrm{T}_{G''/S}(\mathcal{M})\to 1\)
and
\[
1\to\Lie(G'/S,\mathcal{M})\to\Lie(G/S,\mathcal{M})
\to\Lie(G''/S,\mathcal{M})\to 1
\]
are then exact.
\item Let \(1\to H'\to H\to H''\to 1\) be a second sequence of groups; it is
exact if and only if the sequence below is exact:
\[
1\to G'\times_S H'\to G\times_S H\to G''\times_S H''\to 1.
\]
\item If two of the \(S\)-groups \(G'\), \(G\), \(G''\) satisfy~(E), the third
also satisfies~(E).
\item If \(0\to F'\to F\to F''\to 0\) is an exact sequence of
\(\OS\)-modules and if two of the modules \(F'\), \(F\), \(F''\) are good,
the third is also good.
\item If two of the \(S\)-groups \(G'\), \(G\), \(G''\) are good, the third
is also good.
\end{enumeratei}
\end{lemma}

\begin{proof}
\renewcommand{\qedsymbol}{}
{\renewcommand{\thefootnote}{(93)}\nde{One has added the proof of
points~(iii) and~(v).}}
The first part of~(i) is immediate, and the second part follows from it.
Likewise, (ii) is immediate. Let us prove~(iii). For brevity, let us write
\(L(\mathcal{M})=\Lie(G,\mathcal{M})\),
\(L'(\mathcal{M})=\Lie(G',\mathcal{M})\), etc. Then one has the following
commutative diagram
\[
\xymatrix{
1 \ar[r]
  & L'(\mathcal{M}\oplus\mathcal{N}) \ar[r] \ar[d]
  & L(\mathcal{M}\oplus\mathcal{N}) \ar[r] \ar[d]
  & L''(\mathcal{M}\oplus\mathcal{N}) \ar[r] \ar[d]
  & 1
\\
1 \ar[r]
  & L'(\mathcal{M})\times_S L'(\mathcal{N}) \ar[r]
  & L(\mathcal{M})\times_S L(\mathcal{N}) \ar[r]
  & L''(\mathcal{M})\times_S L''(\mathcal{N}) \ar[r]
  & 1
}
\]
in which the first (\resp\ the second) row is exact by~(i) (\resp\ by~(i)
and~(ii)). Assertion~(iii) follows from this.

Let us prove~(iv). One has a commutative diagram
\[
\xymatrix{
0 \ar[r]
  & F'\otimes_{\OS}\mathrm{T}_{\OS/S}(\mathcal{M}) \ar[r] \ar[d]
  & F\otimes_{\OS}\mathrm{T}_{\OS/S}(\mathcal{M}) \ar[r] \ar[d]
  & F''\otimes_{\OS}\mathrm{T}_{\OS/S}(\mathcal{M}) \ar[r] \ar[d]
  & 0
\\
0 \ar[r]
  & \mathrm{T}_{F'/S}(\mathcal{M}) \ar[r]
  & \mathrm{T}_{F/S}(\mathcal{M}) \ar[r]
  & \mathrm{T}_{F''/S}(\mathcal{M}) \ar[r]
  & 0
}
\]
with exact rows (the first, because \(\mathrm{T}_{\OS/S}(\mathcal{M})\) is a
free \(\OS\)-module, hence flat, the second by~(i)). It follows that if two
of the modules \(F'\), \(F\), \(F''\) are good, the third is also good.
Finally, (v) follows from~(iii) and~(iv).
\end{proof}

\begin{lemma}{5.2.2}
\label{II.5.2.2}
Let \(G\) be an \(S\)-group, \(E\), \(H\) two \(G\)-objects, and \(F\) a
\(G\)-\(\OS\)-module.
\begin{enumeratei}
\item The canonical homomorphism
\(E^G\times_S H^G\to(E\times_S H)^G\) is an isomorphism.
\item If \(F\) is a good \(\OS\)-module, then so is \(F^G\).
\end{enumeratei}
\end{lemma}

\begin{proof}
\renewcommand{\qedsymbol}{}
{\renewcommand{\thefootnote}{(94)}\nde{One has set out the original in more
detail in what follows.}}
The first assertion is immediate; let us prove the second. For every
\(S'\to S\), one has a commutative diagram
\[
\xymatrix{
F^G\bigl(I_{S'}(\mathcal{M})\bigr) \ar@{^{(}->}[r]
  & F\bigl(I_{S'}(\mathcal{M})\bigr) \\
F^G(S')\otimes_{\mathrm{O}(S')}\mathrm{O}\bigl(I_{S'}(\mathcal{M})\bigr)
  \ar@{^{(}->}[r] \ar[u]_{\varphi}
  & F(S')\otimes_{\mathrm{O}(S')}\mathrm{O}\bigl(I_{S'}(\mathcal{M})\bigr)
  \ar[u]^{\simeq}_{\varphi_1}
}
\]
and one must prove that \(\varphi\) is bijective; now, it is evidently
injective; let us show that it is surjective. Let \((t_0,\ldots,t_n)\) be a
basis of the free \(\mathrm{O}(S')\)-module
\(\mathrm{O}(I_{S'}(\mathcal{M}))\) and let
\[
u=\sum_{i=0}^{n}f_i\otimes t_i
\]
be an element of
\(F(S')\otimes_{\mathrm{O}(S')}\mathrm{O}(I_{S'}(\mathcal{M}))\) such that
\(\varphi_1(u)\) belongs to \(F^G(I_{S'}(\mathcal{M}))\). Let us show that
the \(f_i\) belong to \(F^G(S')\).

Let \(S''\to S'\); one may regard \(S''\) as lying over
\(I_{S'}(\mathcal{M})\) via the zero section \(\varepsilon_{\mathcal{M}}\).
Then, for every \(g\in G(S'')\), one has:
\[
u_{S''}=g\cdot u_{S''}
=\sum_{i=0}^{n}g\cdot(f_i)_{S''}\otimes(t_i)_{S''}.
\]
As the \((t_i)_{S''}\) form a basis of
\(\mathrm{O}(I_{S''}(\mathcal{M}))
\simeq\mathrm{O}(I_{S'}(\mathcal{M}))\otimes_{\mathrm{O}(S')}\mathrm{O}(S'')\)
over \(\mathrm{O}(S'')\), it follows that
\(g\cdot(f_i)_{S''}=(f_i)_{S''}\), whence \(f_i\in F^G(S')\) for every~\(i\).
\end{proof}

\par\addvspace{0.7\baselineskip}%
\noindent\textbf{Notations.}\ ---
{\renewcommand{\thefootnote}{(95)}\nde{One has set out these notations in
more detail.}}
If \(E\) is an \(S\)-group and \(F\) a sub-\(S\)-group of \(E\), one denotes
by \(E/F\) the \(S\)-functor which to every \(S'\to S\) associates the set
\(E(S')/F(S')\) of classes \(x=x\,F(S')\), \(x\in E(S')\). If \(E\) is a
commutative \(S\)-group then \(E/F\) is a commutative \(S\)-group.

Let now \(G\) be an \(S\)-group and \(K\) a sub-\(S\)-group of \(G\); set
\(E=\Lie(G/S,\mathcal{M})\) and \(F=\Lie(K/S,\mathcal{M})\). The adjoint
action of \(K\) on \(E\) stabilizes \(F\), hence induces an action of \(K\)
on the \(S\)-functor \(E/F\). Then, for every \(S'\to S\), one has:
\[
(E/F)^K(S')
=
\left\{
x\in E(S')/F(S')
\;\middle|\;
\begin{aligned}
&\text{for every }f\colon S''\to S'\text{ and }k\in K(S''),\\
&f^*(x^{-1})\,\Ad(k)\bigl(f^*(x)\bigr)\in F(S'')
\end{aligned}
\right\},
\]
where \(f^*(x)\) denotes the image of \(x\) in \(E(S'')\).

\begin{theorem}{5.2.3}
\label{II.5.2.3}
Let \(G\) be an \(S\)-group, and \(K\) a sub-\(S\)-group of \(G\). Let us denote
(I~2.3.3)
\[
\mathrm{N}=\Norm_G(K),
\qquad
\mathrm{Z}=\Centr_G(K).
\]
Let us make \(K\) operate on \(\Lie(G/S,\mathcal{M})\) by means of the adjoint
representation of \(G\).
\begin{enumeratei}
\item If the group law of \(\Lie(G/S,\mathcal{M})\) is
commutative{\renewcommand{\thefootnote}{(*)}\footnote{Condition
automatically satisfied if \(G\) satisfies~(E) (\cf~3.9), for example if
\(G\) is representable.}}%
{\renewcommand{\thefootnote}{(96)}\nde{For an example where the \(S\)-group
\(\Lie(G/S)\) is not commutative, see~6.3.}},
then
\[
\Lie(N/S,\mathcal{M})\big/\Lie(K/S,\mathcal{M})
=
\bigl(\Lie(G/S,\mathcal{M})\big/\Lie(K/S,\mathcal{M})\bigr)^{K}.
\]
\item If the group law of \(\Lie(G/S,\mathcal{M})\) is commutative~\((*)\),
then
\[
\Lie(Z/S,\mathcal{M})=\Lie(G/S,\mathcal{M})^{K}.
\]
\item If \(G\) satisfies~(E) (\resp\ if \(G\) and \(K\) satisfy~(E)), then
\(Z\) satisfies~(E) (\resp~\(N\) satisfies~(E)).
\item Suppose \(G\) is good; then \(Z\) is good; if moreover \(G\) is very
good, then \(Z\) is very good.
\item Suppose \(G\) and \(K\) are good; then \(N\) is good; if moreover \(G\)
is very good, then \(N\) is very good.
\end{enumeratei}
\end{theorem}

To prove~(i) and~(ii)%
{\renewcommand{\thefootnote}{(97)}\nde{One has set out the proof in more
detail, adding in particular Lemma~5.2.3.0, implicit in the original.}}
we shall use the following lemma, which results from the diagram of exact
sequences considered in~4.1.B (with \(G\) and \(H\) interchanged).

\begin{lemma}{5.2.3.0}
\label{II.5.2.3.0}
Let \(G\) be an \(S\)-group, \(H\) a sub-\(S\)-group of \(G\), and
\(\mathcal{M}\) a free \(\OS\)-module of finite type. Then
\(\mathrm{T}_{H/S}(\mathcal{M})\) and \(\Lie(G/S,\mathcal{M})\) are
sub-\(S\)-groups of \(\mathrm{T}_{G/S}(\mathcal{M})\) and one has:
\[
\mathrm{T}_{H/S}(\mathcal{M})\cap\Lie(G/S,\mathcal{M})
=\Lie(H/S,\mathcal{M}),
\]
where one has set
\(\mathrm{T}_{H/S}(\mathcal{M})\cap\Lie(G/S,\mathcal{M})
\overset{\mathrm{def}}{=}
\mathrm{T}_{H/S}(\mathcal{M})
\times_{\mathrm{T}_{G/S}(\mathcal{M})}
\Lie(G/S,\mathcal{M})\).
\end{lemma}

As the functors considered in~(i) and~(ii) commute with base change, it
suffices to show the equalities of \(S\)-points.

Set \(H=N\) (resp.\ \(=Z\)) and let \(\alpha=\pm 1\). By the preceding lemma
and the definition of \(N\) and \(Z\) (\cf~I~2.3.3), one has:
\(\Lie(H/S,\mathcal{M})(S)=\)
\[
\left\{
\begin{gathered}
X\in\Lie(G/S,\mathcal{M})(S)\subset G\bigl(I_S(\mathcal{M})\bigr)
\;\big|\;\\
\begin{aligned}
&\text{for every }f\colon S'\to I_S(\mathcal{M})\text{ and }u\in K(S'),\\
&f^*(X^{\alpha})\cdot u\cdot f^*(X^{-\alpha})\cdot u^{-1}\in K(S')\\
&\text{resp.\ }f^*(X)\cdot u\cdot f^*(X^{-1})\cdot u^{-1}=1
\end{aligned}
\end{gathered}
\right\}
\quad(*),\qquad
\]
where \(f^*(X)\) denotes the image of \(X\) in \(G(S')\).

To simplify the writing, let us set
\[
\mathfrak{g}=\Lie(G/S,\mathcal{M}),\quad
\mathfrak{k}=\Lie(K/S,\mathcal{M}),\quad
\mathfrak{n}=\Lie(N/S,\mathcal{M}),\quad
\mathfrak{z}=\Lie(Z/S,\mathcal{M}).
\]
If \(X\in\Lie(H/S,\mathcal{M})(S)\), the equalities \((*)\) are valid for
every \(f\colon S'\to S\), since \(f=\rho\circ\varepsilon_{\mathcal{M}}\circ f\)
factors through \(I_S(\mathcal{M})\) (where
\(\rho\colon I_S(\mathcal{M})\to S\) is the structural morphism and
\(\varepsilon_{\mathcal{M}}\colon S\to I_S(\mathcal{M})\) the zero section).
One deduces from this that
\[
\mathfrak{n}(S)/\mathfrak{k}(S)
\subset(\mathfrak{g}/\mathfrak{k})^{K}(S)
\qquad\text{and}\qquad
\mathfrak{z}(S)\subset\mathfrak{g}^{K}(S).
\]
